\chapter{Training}
\label{ch:training}

\section{The short version}

\begin{lstlisting}
poraque-train --config configs/train.yaml
\end{lstlisting}

This trains \textbf{one} \opt{ext2chg} model and \textbf{one} \opt{chg2tau}
model on the combined data of every structure, and writes:

Everything it writes lands in \emph{one directory named after the run}, so a
trained model and the evidence for it stay together:

\begin{center}
\begin{tabular}{@{}ll@{}}
\toprule
Path & Contents \\
\midrule
\file{models/<name>/<name>.poraque} & both trained operators, in one file \\
\file{models/<name>/log/} & training log, metrics, resolved configuration \\
\file{models/<name>/plots/} & loss curves, cross-sections, parity plots \\
\file{models/<name>/report/} & the typeset PDF report \\
\bottomrule
\end{tabular}
\end{center}

\subsection{What the report's performance table says}

One row per split --- \opt{train}, \opt{validation}, \opt{all} --- and five
columns:

\begin{center}
\begin{tabular}{@{}lp{8.4cm}@{}}
\toprule
Column & What it answers \\
\midrule
rel.\ $L^2$ & The headline error, and the one every other number in this
manual is quoted in. \\
rel.\ $H^1$ & Values \emph{and} gradients. A prediction can match a density
pointwise and still be rough, and $\tauvW$ depends on $\nabla\rho$ --- so a
small $L^2$ beside a large $H^1$ is a model whose energetics will disappoint. \\
MAE & The typical voxel, in the field's own units, where a squared error is not
readable as a quantity. \\
Max.\ error & The worst voxel. Means and RMS both hide a single catastrophic
site, and near a nucleus is exactly where one occurs. \\
$|\Delta N|$ & Conservation. For \opt{ext2chg} this is the electron-count error
in electrons, and it is the metric a pointwise loss controls worst. \\
\bottomrule
\end{tabular}
\end{center}

\begin{ptip}
The \opt{train} and \opt{validation} rows read against each other \emph{are}
the generalisation gap. That is the single most useful thing the table says,
and it is why the splits are rows.
\end{ptip}

\begin{pnote}
The rel.\ $H^1$ here is the textbook relative Sobolev norm, which has no free
parameter --- \textbf{not} the $H^1$ objective \opt{loss: relative\_h1}
minimises, whose value depends on \opt{sobolev\_weight} and so cannot be
compared between two runs. Every column is computed without consulting the
loss, for exactly that reason.

For a two-channel model every number is the \textbf{density} channel; the
magnetisation is measured separately, as \opt{magnetisation\_relative\_l2} in
the metrics JSON, and the report says so in its caveats.

The table lists \emph{splits}, not structures: it used to print one row per
material, so a 115-material run made six pages of numbers nobody reads
individually. The per-structure figures --- and \opt{mse}, \opt{rmse},
\opt{r2}, \opt{nrmse\_range} and \opt{jsd}, which are not typeset --- are all
in \file{log/<name>.json}, which is machine-readable and has no margins.
\end{pnote}

\noindent
\opt{<name>} is \opt{task.name} from the configuration, and it is the only
thing you set: two runs with different names cannot overwrite each other, and
there are no output paths to keep in step by hand. A trained model is not one
file --- it is weights, the numbers that say how good they are, the figures
behind those numbers, and the configuration that produced them --- so they
arrive and leave together.

\section{One model, all structures}

A single set of weights is fitted across every training structure at once, and
batches mix materials --- never one model per material.

\begin{pnote}
There is \textbf{one} training protocol and \textbf{one} variation on it. By
default a fifth of the structures (\opt{valid\_fraction}) is held back for
validation, so the printed score is genuinely held out and
\opt{early\_stopping} has something to watch. Nothing else needs selecting.
\end{pnote}

\begin{pwarn}
Set \opt{valid\_fraction: 0} and nothing is held out, so the metrics become a
\textbf{training fit}: they show the model can represent the data, not how it
will do on a new material. On the reference data the training fit is about
$4\times$ better than the cross-validated score, which is the size of the
mistake if the two are confused.

That setting is still the right one for the artefact you ship, since it uses
all the data --- just quote a cross-validated number alongside it.
\end{pwarn}

\section{Measuring generalisation}

Splitting is always at the \emph{structure} level.

\subsection{The default split}

\begin{lstlisting}
poraque-train --config configs/train.yaml \
       --valid-fraction 0.2
\end{lstlisting}

Structures are shuffled with \opt{seed} and that share is kept back.

\subsection{K-fold cross-validation}
\label{sec:kfold}

This is the only variation on the protocol; \opt{valid\_fraction} is ignored,
since the folds define the splits.

\begin{lstlisting}
poraque-train --config configs/train.yaml \
       --kfold --k-folds 5
\end{lstlisting}

Each fold trains a fresh model on $K-1$ groups and scores it on the held-out
group. A consolidated PDF reports the mean and standard deviation of each
metric across folds.

\begin{ptip}
\opt{k\_folds} is capped at the number of structures; setting it equal to that
count is leave-one-out. Whole materials are held out --- never a subset of
voxels from a material that also appears in training, which would score the
model on data it had effectively already seen.
\end{ptip}

\section{Reading the output}

\begin{center}
\begin{tabular}{@{}ll@{}}
\toprule
Metric & Read it as \\
\midrule
relative $L^2$ & fractional error; $0.03$ means $3\,\%$ \\
$R^2$ & $1$ is perfect; \emph{negative} means worse than the mean \\
MAE, RMSE & absolute error, in the field's own units \\
integral error & how well the electron count or $T_s$ is reproduced \\
\bottomrule
\end{tabular}
\end{center}

For \opt{chg2tau} the log also prints the classical orbital-free functionals
(Thomas--Fermi, von Weizsäcker) evaluated on the same input. Beating those is
the bar --- a learned functional that does not is not adding anything.

\section{Reference numbers}

Seventeen platinum supercells --- ten 27-atom and seven 32-atom, four grid shapes
--- at $32^3$ working resolution, 300 epochs with early stopping, a single
80/20 split at \opt{seed=42}:

\begin{center}
\begin{tabular}{@{}lrr@{}}
\toprule
Model & held out (3 structures) & all structures (training fit) \\
\midrule
\opt{ext2chg} & $0.0379 \pm 0.0027$ & $0.0209$ \\
\opt{chg2tau} & $0.0511 \pm 0.0031$ & $0.0140$ \\
\bottomrule
\end{tabular}
\end{center}

Quote the held-out column. The gap to the training fit is the memorisation
margin --- roughly two- to four-fold here.

\begin{pwarn}
\textbf{Read the held-out column knowing what is in it.} At
\opt{valid\_fraction: 0.2} with 17 structures the validation set is three
structures, and at \opt{seed=42} all three are 32-atom cells. So these numbers
describe the harder cell size only, there is no held-out 27-atom measurement in
them, and the $\pm$ is the spread over three structures rather than a
cross-validation error bar.

Cross-validation on the earlier 12-structure dataset measured the cell-size
effect directly: $0.0189 \pm 0.0047$ (\opt{ext2chg}) on 27-atom cells against
$0.0441 \pm 0.0180$ on 32-atom ones, so a held-out 32-atom cell was about
$2.4\times$ harder. The current single-split figure for that harder subset,
$0.0379$, is better than the $0.0441$ it scored then --- which is what having
four 32-atom cells in training rather than one should buy.

Expect a similar penalty whenever a new cell size enters the dataset, and use
\opt{--kfold} (Section~\ref{sec:kfold}) when you need a figure that covers
every structure rather than whichever three the seed selected.
\end{pwarn}

\section{Options worth knowing}

\begin{description}
  \item[\opt{--pauli-head}] For \opt{chg2tau}, predicts
    $\tau = \tau_{\rm vW}[\rho] + s\,\mathrm{softplus}(f)$ so the exact bound
    $\tau \ge \tau_{\rm vW}$ holds by construction. Improves accuracy and
    trains slightly faster; recommended.
  \item[\opt{--gaussian-blur}] Smooths the external potential. Measured to make
    no meaningful difference at $32^3$ and to remove information near the ionic
    cores; leave it off unless working at higher resolution.
  \item[\opt{--device}] \opt{auto}, \opt{cuda}, \opt{mps} or \opt{cpu}. An
    unavailable backend warns and falls back rather than failing.
  \item[\opt{--resolution}] Working grid. Higher is more faithful and much
    slower; $32$ is a reasonable starting point.
\end{description}
